\chapter{Options de taux vanille}\label{ch:rc:vanilla-rates-options}

En juin 2014, la BCE fixe son taux de facilité de dépôt à $-0.10~\%$, et en moins de deux ans
les taux courts en euros et de nombreux taux forward passent sous zéro. Le modèle de Black, dans
lequel un taux est lognormal, ne permet pas d'évaluer une option sur un forward négatif, ou
sur un strike négatif~: sa formule prend le logarithme des deux. Pendant plus
de vingt ans, les caps et les swaptions étaient cotés en volatilités de Black~; les
cotations qui ne pouvaient exister commencent à apparaître en volatilités normales, en points de base
par an, ou en volatilités d'un taux décalé. Les prix se moquaient du modèle
servant à les coter~; les chiffres de risque, eux, en dépendaient. Ce chapitre évalue
les deux options de taux vanille, caplets et swaptions, dans les deux modèles, réalise le
\omterm{def:rc:vanilla-rates-options:strip}{stripping des caplets} à partir des caps cotés, et montre ce que les modalités de
règlement-livraison d'une swaption pèsent sur son prix. Les produits et leur
marché sont ceux du livre~2, chapitre~13~; les courbes en euros sont celles
du chapitre~2.

\section{Black et Bachelier pour les taux}

Un caplet sur l'Euribor à six mois, fixé en $T_{k-1}$ et payé en $T_k$, verse
$N\delta_k\max(F_k(T_{k-1})-K,0)$, où $F_k(t)$ est le forward de ce fixing.
Sous la probabilité forward de $T_k$, $\mathbb Q^{T_k}$, de la \omterm{def:rc:multi-curve-and-collateral-discounting:curves}{courbe d'actualisation} (livre~4,
chapitre~5), $F_k$ est une martingale, de même que le taux au pair
$S_{a,b}$ d'un swap sous la probabilité annuité $\mathbb Q^{a,b}$. Les deux options sont
donc des options sur une martingale, et le modèle n'a qu'à dire comment cette
martingale est distribuée.

\begin{proposition}[Caplets et swaptions sous leurs probabilités naturelles]\label{prop:rc:vanilla-rates-options:measures}
\[
\mathrm{Caplet}_0 = N\delta_kP(0,T_k)\,\E^{T_k}\bigl[(F_k(T_{k-1})-K)^+\bigr],
\qquad
\mathrm{Payeur}_0 = N A_{a,b}(0)\,\E^{a,b}\bigl[(S_{a,b}(T_a)-K)^+\bigr].
\]
Si $dF = \sigma_N\,dW$, l'espérance est donnée par la formule de Bachelier (livre~2,
chapitre~13)~; si $dF = \sigma F\,dW$, c'est la formule de Black (livre~5,
chapitre~3), avec le forward à la place du spot et sans dérive.
\end{proposition}
\begin{proof}[Preuve partielle]
Prenons $P(t,T_k)$ comme numéraire~: le prix du caplet divisé par ce numéraire est une martingale,
et en $T_k$ le numéraire vaut un. $F_k(t) = (P(t,T_{k-1})/P(t,T_k)-1)/\delta_k$
est un actif négocié divisé par le numéraire, donc une martingale. Le même
raisonnement avec l'annuité comme numéraire donne la swaption~; avec plusieurs courbes,
le forward projeté est défini comme l'espérance sous $\mathbb Q^{T_k}$ du
fixing.
\end{proof}

\begin{definition}[Volatilité lognormale]\label{def:rc:vanilla-rates-options:lognormal}
La \emph{volatilité lognormale}\index{volatilité lognormale} $\sigma_B$ d'une
option de taux est la volatilité pour laquelle la formule de Black, appliquée au
taux forward, reproduit son prix~; c'est une volatilité relative, un pourcentage
du taux, et n'existe que pour un forward et un strike positifs.
\end{definition}

\begin{proposition}[Volatilités normales et lognormales à la monnaie]\label{prop:rc:vanilla-rates-options:atm}
Pour une option à la monnaie ($K=F>0$), les deux modèles donnent le même prix lorsque
\[
\sigma_N\sqrt{T}\,\frac{1}{\sqrt{2\pi}} = F\bigl(2\Phi(\tfrac12\sigma_B\sqrt T)-1\bigr),
\quad\text{soit}\quad
\sigma_N \approx \sigma_B F \ \text{ pour de petits }\sigma_B\sqrt T.
\]
\end{proposition}
\begin{proof}
Les deux prix à la monnaie sont $\sigma_N\sqrt{T}/\sqrt{2\pi}$ et
$F(2\Phi(\sigma_B\sqrt T/2)-1)$~; développer $\Phi$ au premier ordre.
\end{proof}

Un modèle normal évalué avec une seule volatilité pour tous les strikes présente un skew négatif
marqué en termes lognormaux (la \cref{fig:rc:vanilla-rates-options:smile})~: pour une
option à cinq ans sur un forward de 2.50~\%, à 80 points de base de
volatilité normale, la volatilité de Black vaut 32.7~\% à la monnaie (l'approximation
donne $80/250 = 32~\%$), 71.1~\% pour un strike de 0.50~\% et 23.8~\% à 4.50~\%. Un
marché dont le smile est plat en termes normaux paraît donc fortement incliné pour
quiconque cote en termes de Black, et le choix d'un modèle est aussi un choix de la façon dont le
smile évolue quand les taux évoluent.

\begin{omfigure}
\begin{tikzpicture}
\begin{axis}[width=11cm,height=5.4cm,
  xlabel={strike (\%)},ylabel={volatilité lognormale (\%)},
  tick label style={font=\small},label style={font=\small},
  xmin=0.5,xmax=5,ymin=15,ymax=75,grid=major,grid style={gray!25},
  table/col sep=comma]
\addplot[omDef,very thick] table[x=k,y=black]{figdata/rates-credit-risk/04-vanilla-rates-options/smile.csv};
\draw[gray,dashed] (axis cs:2.5,15) -- (axis cs:2.5,75) node[pos=0.9,right,font=\footnotesize,text=black] {forward 2.50~\%};
\end{axis}
\end{tikzpicture}

\omcaption{La volatilité implicite lognormale (de Black) des options à cinq ans sur un
forward de 2.50~\% évaluées avec une seule volatilité normale de 80 points de base.
Un smile normal plat est un skew lognormal marqué. Données~: le
tutoriel du chapitre.}\label{fig:rc:vanilla-rates-options:smile}
\end{omfigure}

\begin{dated}{2026-09}{Taux négatifs et conventions des options de taux}\label{dat:rc:vanilla-rates-options:negative}
Le taux de la facilité de dépôt de la BCE a été négatif du 11 juin 2014 ($-0.10~\%$) au 27
juillet 2022, avec un point bas à $-0.50~\%$ depuis le 18 septembre 2019. Les options sur
forwards ou strikes négatifs n'ont pas de volatilité de Black, et le marché a commencé à coter les caps,
floors et swaptions aussi en volatilités normales ou en volatilités de Black décalées
(chapitre~5). Depuis le 26 novembre 2018, la méthode de règlement en espèces par défaut pour
les swaptions en euros dans la matrice de règlement-livraison de l'ISDA est devenue le prix cash
collatéralisé, au lieu de la méthode du taux au pair.
\end{dated}

\section{Caps, floors et stripping des caplets}

Un cap est coté par un seul nombre, son prix ou une seule volatilité pour tous ses
caplets. Les caplets d'un cap à cinq ans sont les mêmes que les quatre premières
années d'un cap à dix ans, et ne peuvent pas avoir deux prix.

\begin{definition}[Volatilité plate, volatilité de caplet]\label{def:rc:vanilla-rates-options:flat}
La \emph{volatilité plate}\index{volatilité plate} d'un cap est la volatilité unique
pour laquelle la somme des prix de ses caplets égale le prix du cap. Une
\emph{volatilité de caplet}\index{volatilité de caplet} est la volatilité d'un
caplet, fonction de sa date de fixing (et du strike)~; les volatilités de caplet de
tous les caps d'un même strike forment une seule structure par terme.
\end{definition}

\begin{definition}[Stripping des caplets]\label{def:rc:vanilla-rates-options:strip}
Le \emph{stripping des caplets}\index{stripping des caplets} est la calibration
d'une structure par terme de \omterm{def:rc:vanilla-rates-options:flat}{volatilités de caplet} qui reproduit le prix de chaque cap coté d'un même strike~:
avec des caps de maturité croissante, les caplets compris entre deux maturités consécutives
reçoivent une seule volatilité, choisie pour que le cap le plus long reproduise son prix,
les \omterm{def:rc:vanilla-rates-options:flat}{volatilités des caplets} plus courts restant inchangées.
\end{definition}

\omcode{code/firm/capfloor/firm_capfloor.py}{63}{86}{Stripping des volatilités de caplet constantes par morceaux à partir des volatilités plates des caps.}\label{lst:rc:vanilla-rates-options:strip}

\begin{example}[Un stripping de caps en euros]\label{ex:rc:vanilla-rates-options:strip}
Volatilités normales plates des caps de strike 3~\% sur l'Euribor six mois de 58, 66,
72, 75, 76, 75 et 72 points de base pour des maturités d'un, deux, trois, quatre, cinq, sept et
dix ans (à titre illustratif~; le premier caplet, déjà fixé, exclu) se décomposent en
\omterm{def:rc:vanilla-rates-options:flat}{volatilités de caplet} de 58.0, 66.6, 75.7, 78.9, 78.0, 73.7 et 67.9 points de base
sur les intervalles successifs (\cref{fig:rc:vanilla-rates-options:strip}). Là où
les \omterm{def:rc:vanilla-rates-options:flat}{volatilités plates} montent, les nouveaux caplets doivent se situer au-dessus du niveau
plat pour relever la moyenne~; la bosse de la courbe des caplets est plus haute et plus tôt que celle de la
courbe plate.
\end{example}

\begin{omfigure}
\begin{tikzpicture}
\begin{axis}[width=12cm,height=5.4cm,
  xlabel={date de fixing du caplet (années)},ylabel={volatilité normale (pb)},
  tick label style={font=\small},label style={font=\small},
  xmin=0,xmax=10,ymin=50,ymax=85,grid=major,grid style={gray!25},
  legend columns=2,legend style={font=\footnotesize,at={(0.5,-0.25)},anchor=north,draw=none,fill=none,/tikz/every even column/.append style={column sep=8pt}},
  table/col sep=comma]
\addplot[omThm,very thick,const plot,mark=*,mark size=1.3pt] table[x=t,y=caplet]{figdata/rates-credit-risk/04-vanilla-rates-options/strip.csv};
\addplot[omDef,thick,dashed,const plot] table[x=t,y=flat]{figdata/rates-credit-risk/04-vanilla-rates-options/strip.csv};
\legend{volatilité de caplet issue du stripping,volatilité plate du cap le plus court contenant le caplet}
\end{axis}
\end{tikzpicture}

\omcaption{\omterm{def:rc:vanilla-rates-options:flat}{Volatilités de caplet} obtenues par stripping à partir des \omterm{def:rc:vanilla-rates-options:flat}{volatilités plates} des caps
de strike 3~\%, une par intervalle entre les maturités cotées. Données~: les
cotations illustratives sur les courbes en euros du chapitre~2~; le
tutoriel du chapitre.}\label{fig:rc:vanilla-rates-options:strip}
\end{omfigure}

\begin{omfigure}
\begin{tikzpicture}[x=1.15cm,y=1cm,font=\small]
\draw[->,thick] (0,0) -- (10.6,0) node[right] {temps};
\foreach \x/\l in {0/spot,1/$T_0$,2/$T_1$,3/$T_2$,4/$T_3$,9/$T_{n-1}$,10/$T_n$} {\draw (\x,-0.1) -- (\x,0.1); \node[below,font=\footnotesize] at (\x,-0.12) {\l};}
\node at (6.5,0.35) {$\dots$};
\foreach \a/\b/\y in {1/2/1.1,2/3/1.1,3/4/1.1,9/10/1.1} {\draw[omDef,thick,fill=omDef!12] (\a,\y) rectangle (\b,\y+0.45);}
\node[font=\footnotesize] at (1.5,1.32) {1};
\node[font=\footnotesize] at (2.5,1.32) {2};
\node[font=\footnotesize] at (3.5,1.32) {3};
\node[font=\footnotesize] at (9.5,1.32) {$n$};
\node[font=\footnotesize,anchor=west] at (4.3,1.32) {caplets};
\draw[dashed,gray,thick] (0,1.1) rectangle (1,1.55); \node[font=\footnotesize,gray] at (0.5,1.32) {fixé};
\draw[-{Stealth[length=2mm]},omThm] (2,1.05) -- (2,0.15);
\draw[-{Stealth[length=2mm]},omExo] (3,1.05) -- (3,0.15);
\node[font=\footnotesize,text=omThm,anchor=east] at (1.8,2.3) {caplet 2 se fixe à $T_1$};
\node[font=\footnotesize,text=omExo,anchor=west] at (3.3,2.3) {et paie à $T_2$};
\draw[omThm,thin] (1.8,2.3) -- (2,1.6);
\draw[omExo,thin] (3.3,2.3) -- (3,1.6);
\end{tikzpicture}

\omcaption{Un cap sur l'Euribor six mois à partir du spot~: le taux de chaque caplet se fixe
au début de sa période et est payé à la fin. Le taux de la première période est
déjà fixé à la négociation, de sorte que les caps du marché commencent au deuxième caplet.}\label{fig:rc:vanilla-rates-options:timeline}
\end{omfigure}

\section{Swaptions~: règlement-livraison physique et règlement en espèces}

Une swaption à règlement-livraison physique livre le swap~; sa valeur est l'annuité
multipliée par une option sur le taux swap forward
(\cref{prop:rc:vanilla-rates-options:measures}). En euros et en livres sterling, la plupart des
swaptions se règlent en espèces, et le montant de règlement exige une annuité que
les deux parties puissent calculer à la date d'exercice à partir d'un seul nombre.

\begin{definition}[Swaption à règlement en espèces, annuité cash]\label{def:rc:vanilla-rates-options:cash}
Une \emph{swaption à règlement en espèces}\index{swaption à règlement en espèces} verse, à
l'exercice, la valeur du swap sous-jacent calculée selon une méthode convenue
plutôt que de la livrer. Selon la méthode du taux au pair, la valeur est
$N\,a(S_T)\,(S_T-K)^+$ avec l'\emph{annuité cash}\index{annuité cash}
\[
a(S) = \sum_{i=1}^{nm}\frac{1/m}{(1+S/m)^i},
\]
l'annuité du swap actualisée à son propre taux au pair $S$ ($m$ paiements par
an~)~; selon la méthode du prix cash collatéralisé, elle est la valeur du swap sur la
\omterm{def:rc:multi-curve-and-collateral-discounting:curves}{courbe d'actualisation} de la chambre de compensation convenue, comme en règlement-livraison physique.
\end{definition}

\omcode{code/firm/capfloor/firm_capfloor.py}{110}{121}{L'annuité cash et la formule traditionnelle de règlement en espèces au taux au pair.}\label{lst:rc:vanilla-rates-options:cash}

\begin{example}[Cinq ans sur dix ans]\label{ex:rc:vanilla-rates-options:swaption}
Sur les courbes du chapitre~2, le swap à dix ans débutant dans cinq ans a un taux au pair
forward de 3.098~\% contre l'Euribor six mois et une annuité (calculée sur la \omterm{def:rc:multi-curve-and-collateral-discounting:curves}{courbe d'actualisation}) de
7.740. À 80 points de base de volatilité normale (une \omterm{def:rc:vanilla-rates-options:lognormal}{volatilité lognormale} de
26.2~\%), le payeur physique à la monnaie sur 100~millions d'euros vaut
5\,525\,150~euros. La formule cash traditionnelle, $P(0,T)\,a(S_0)\times$ option,
donne 5\,410\,452~euros, soit 2.1~\% de moins~: $P(0,T)\,a(S_0) = 7.579$, car
l'\omterm{def:rc:vanilla-rates-options:cash}{annuité cash} s'actualise au taux du swap Euribor, au-dessus des taux €STR de
l'annuité collatéralisée. Le payoff au taux au pair n'est pas non plus une fonction du
numéraire de la probabilité annuité, de sorte que la formule est elle-même une approximation~; le
prix cash collatéralisé lève les deux problèmes en se réglant à la
valeur physique.
\end{example}

\section{Vega, le smile et le choix du modèle}

Le vega normal d'un caplet, $\partial\mathrm{Caplet}/\partial\sigma_N =
N\delta P\sqrt T\varphi(d)$ avec $d=(F-K)/(\sigma_N\sqrt T)$, est maximal à la
monnaie et ne dépend pas du niveau des taux~; le vega Black est
proportionnel à $F$. Un desk qui couvre dans un modèle et est valorisé dans l'autre
voit son vega varier quand les taux bougent sans que le marché ait changé. Deux règles en découlent.
Couvrir dans le modèle où le smile du marché est le plus stable (pour les taux euro et
dollar, modèles normaux ou décalés depuis les années de taux négatifs)~; et
rapporter le vega par point de base de volatilité normale, une mesure additive entre
caplets et swaptions de forwards différents.

\begin{omfigure}
\begin{tikzpicture}
\begin{axis}[width=11cm,height=5.2cm,
  xlabel={strike du cap (\%)},ylabel={prime (pb par an)},
  tick label style={font=\small},label style={font=\small},
  xmin=2,xmax=4.5,ymin=0,ymax=75,grid=major,grid style={gray!25},
  table/col sep=comma]
\addplot[omDef,very thick,mark=*,mark size=1.4pt] table[x=k,y=bp]{figdata/rates-credit-risk/04-vanilla-rates-options/capstrike.csv};
\draw[gray,dashed] (axis cs:3,0) -- (axis cs:3,75) node[pos=0.88,right,font=\footnotesize,text=black] {le 3~\% du trésorier};
\end{axis}
\end{tikzpicture}

\omcaption{Coût courant d'un cap à cinq ans sur l'Euribor six mois, sur les courbes du chapitre~2
par strike, avec les \omterm{def:rc:vanilla-rates-options:flat}{volatilités de caplets} strippées à 3~\% utilisées à chaque
strike (le smile ignoré). Le coût passe de 70 points de base par an à un
strike de 2~\% jusqu'à 4 points de base à 4.5~\%. Données~: le tutoriel du chapitre.}\label{fig:rc:vanilla-rates-options:capstrike}
\end{omfigure}

\begin{remark}[Plat ou strippé~: même prix, risque différent]\label{rem:rc:vanilla-rates-options:samerisk}
Le cap à cinq ans du problème ci-dessous a le même prix à sa \omterm{def:rc:vanilla-rates-options:flat}{volatilité plate}
et aux \omterm{def:rc:vanilla-rates-options:flat}{volatilités de caplets} strippées, par construction~; son vega total
est lui aussi le même, mais son vega par compartiment ne l'est pas~: la courbe strippée le place
sur les échéances propres des caplets, là où se trouvent les caps et swaptions de couverture.
\end{remark}

\section{Tutoriel~: stripper, valoriser, comparer}\label{tut:rc:vanilla-rates-options:tutorial}

\begin{tutorial}
\textbf{Objectif.} Stripper les \omterm{def:rc:vanilla-rates-options:flat}{volatilités de caplets} à partir des cotations de caps, valoriser un cap et une
swaption sur les courbes en euros du chapitre~2, et comparer les méthodes de règlement.
\textbf{État final~:}
les \cref{fig:rc:vanilla-rates-options:smile,fig:rc:vanilla-rates-options:strip} et les
valeurs de l'\cref{ex:rc:vanilla-rates-options:swaption}.
\begin{tutsteps}
\item \textbf{Smile}~: \texttt{smile\_from\_flat\_normal()} convertit une volatilité normale plate
en volatilités de Black par strike, à l'aide de
\texttt{normal\_to\_black} du livre~2.
\item \textbf{Strip}~: \texttt{stripped()} renvoie sept nœuds~; vérifier que chaque
cap se revalorise à sa \omterm{def:rc:vanilla-rates-options:flat}{volatilité plate}.
\item \textbf{Swaption}~: \texttt{swaption\_table()} donne le forward, les deux
annuités et les deux prix.
\item \textbf{Cap}~: \texttt{treasurer\_cap()} et \texttt{cap\_strike\_table()}~;
\texttt{fig\_rc\_vanilla.py} produit les graphiques.
\end{tutsteps}
\textbf{Pour aller plus loin.} Valoriser le cap avec une \omterm{def:rc:vanilla-rates-options:lognormal}{volatilité lognormale} plate de
25~\% et comparer~; stripper avec des caps cotés chaque année au lieu d'un, deux,
trois, quatre, cinq, sept et dix ans.
\end{tutorial}

\section{Construction~: caps, floors et swaptions}\label{bld:rc:vanilla-rates-options:capfloor}

\begin{build}
\bfield{Rôle} Les options de taux vanille de la petite entreprise fictive, valorisées sur la
multi-courbe~; les instruments de calibration des modèles des chapitres 5 à 10, et
les options que le moteur de risque du \cref{ch:rc:build-a-risk-engine} enregistre
auprès de la bibliothèque de valorisation.
\bfield{Interface} \texttt{Cap(start, years, strike, months, floor, skip\_first)}~;
\texttt{caplets}, \texttt{cap\_price(cap, proj, disc, vol, model)},
\texttt{cap\_vega}, \texttt{strip\_caplet\_vols}, \texttt{piecewise}~;
\texttt{forward\_swap}, \texttt{swaption\_physical}, \texttt{cash\_annuity},
\texttt{swaption\_cash\_irr}.
\bfield{Règles} Caplets prospectifs fixés au début de leur période~; le
premier caplet ignoré par défaut~; \texttt{model} vaut \texttt{normal} ou
\texttt{lognormal}~; le \texttt{firm\_normalvol} du livre~2 est importé, non modifié.
\bfield{Tests d'acceptation} \texttt{code/firm/capfloor/tests/}~: une volatilité nulle
donne la valeur intrinsèque~; cap moins floor égale le swap à terme~; les volatilités strippées
revalorisent chaque cap~; payeur moins receveur égale le swap à terme~; l'\omterm{def:rc:vanilla-rates-options:cash}{annuité cash} a la
forme fermée et baisse avec le taux.
\bfield{Pour aller plus loin} Strip dépendant du strike (un smile de caplets par échéance)~; le
prix cash collatéralisé avec une courbe de règlement distincte~; les caplets rétrospectifs
(chapitre~10).
\end{build}

\begin{omsources}
\item F.\ Black, «~The pricing of commodity contracts~», \emph{Journal of
Financial Economics} 3, 1976.
\item ISDA, «~Market practice change for settlement of EUR swaptions to
collateralized cash price~», novembre 2018.
\item D.\ Brigo et F.\ Mercurio, \emph{Interest Rate Models: Theory and
Practice}, 2e éd., Springer, 2006, chapitre 6.
\item BCE, \emph{Key ECB interest rates}.
\end{omsources}

\section{Exercices}

\begin{exercise}[$\star$]\label{exo:rc:vanilla-rates-options:1}
Un caplet sur 10~millions d'euros, de strike 3~\%, fixe à 3.40~\% pour une période
dont la fraction vaut 0.5069. Calculer le montant versé et la date du paiement.
\end{exercise}

\begin{exercise}[$\star$]\label{exo:rc:vanilla-rates-options:2}
Une option à cinq ans à la monnaie sur un forward de 2.50~\% a une volatilité normale
de 80 points de base. Estimer sa \omterm{def:rc:vanilla-rates-options:lognormal}{volatilité lognormale}, et la comparer à la valeur exacte
de 32.7~\%.
\end{exercise}

\begin{exercise}[$\star$]\label{exo:rc:vanilla-rates-options:3}
Expliquer pourquoi les caps de marché sur l'Euribor six mois excluent le premier caplet.
\end{exercise}

\begin{exercise}[$\star\star$]\label{exo:rc:vanilla-rates-options:4}
Le cap à un an a une \omterm{def:rc:vanilla-rates-options:flat}{volatilité plate} de 58 et le cap à deux ans de 66
points de base. Sans calculer, dire si la volatilité strippée des caplets de la deuxième année
est supérieure ou inférieure à 66~? Vérifier avec l'\cref{ex:rc:vanilla-rates-options:strip}.
\end{exercise}

\begin{exercise}[$\star\star$]\label{exo:rc:vanilla-rates-options:5}
Expliquer l'écart de 2.1~\% entre la formule de prix cash au taux au pair et le prix physique
dans l'\cref{ex:rc:vanilla-rates-options:swaption}, et pourquoi le prix cash collatéralisé
referme l'écart.
\end{exercise}

\begin{exercise}[$\star\star$]\label{exo:rc:vanilla-rates-options:6}
Lire la \cref{fig:rc:vanilla-rates-options:smile}~: pourquoi un smile normal plat
ressemble-t-il à un skew lognormal négatif~? Quel smile aurait un marché lognormal
plat, exprimé en termes normaux~?
\end{exercise}

\begin{exercise}[$\star\star\star$]\label{exo:rc:vanilla-rates-options:7}
\emph{Programmation.} Trouver le strike du floor qui rend un collar à cinq ans
(long le cap à 3~\%, short le floor) gratuit, avec les volatilités strippées
à chaque strike.
\end{exercise}

\begin{exercise}[$\star\star\star$]\label{exo:rc:vanilla-rates-options:8}
\emph{Trouver la faille.} «~Le cap à cinq ans est coté 76 points de base, donc je
price les deux caplets qui fixent après quatre ans à 76.~» Qu'est-ce qui ne va
pas, et de combien ce calcul fausse-t-il le prix de ces caplets sur 200~millions d'euros~?
\end{exercise}

\section{Problème~: le cap du trésorier}

\begin{problem}[{Problème du week-end --- plafonner un prêt à taux variable}]\label{pb:rc:vanilla-rates-options:1}
Une entreprise a emprunté 200~millions d'euros pour cinq ans à l'Euribor six mois
majoré d'une marge, le premier taux étant déjà fixé. Son conseil d'administration
veut plafonner le taux à 3~\%. Sur les courbes du chapitre~2 et les volatilités
strippées de l'\cref{ex:rc:vanilla-rates-options:strip}, elle achète le cap à cinq ans à 3~\%.

\medskip\noindent\textbf{Partie I --- Le cap.}
\begin{enumerate}
\item Combien de caplets détient-elle, et pourquoi ce nombre~?
\item Donner sa prime en euros et en pourcentage du notionnel.
\item Exprimer la prime comme un coût annuel en points de base.
\item Le pricer à la \omterm{def:rc:vanilla-rates-options:flat}{volatilité plate} à cinq ans, à la place, et expliquer le résultat.
\item La moyenne des forwards des caplets vaut 2.44~\%. Pourquoi payer quoi que ce soit
pour un cap dont le strike est au-dessus~?
\end{enumerate}

\medskip\noindent\textbf{Partie II --- Risque.}
\begin{enumerate}[resume]
\item Donner son vega par point de base de volatilité normale.
\item Si les volatilités normales augmentent de 10 points de base, de combien le cap gagne-t-il~?
\item Son vega lognormal croît-il ou décroît-il si les taux montent, à volatilité normale
constante~?
\item Qui est short de ce vega, et comment se couvre-t-il~?
\item Pricer le cap à une \omterm{def:rc:vanilla-rates-options:lognormal}{volatilité lognormale} plate de 25~\% et commenter.
\end{enumerate}

\medskip\noindent\textbf{Partie III --- Scénarios.}
\begin{enumerate}[resume]
\item Si chaque fixing restant vaut 3.50~\%, que verse le cap au total,
non actualisé~?
\item Donner la valeur actuelle de ces paiements aujourd'hui.
\item À quel fixing moyen l'entreprise atteint-elle le point mort sur la prime,
à peu près~?
\item Si tous les fixings restent sous 3~\%, qu'a-t-elle acheté~?
\item Pourquoi le cap n'est-il pas une option sur le taux moyen~?
\end{enumerate}

\medskip\noindent\textbf{Partie IV --- Jugement.}
\begin{enumerate}[resume]
\item Comparer le cap avec le fait de payer le taux fixe d'un swap à cinq ans.
\item Donner le strike du floor d'un collar à coût nul et dire ce que l'entreprise
abandonne.
\item Pourquoi une banque pourrait-elle coter ce cap au-dessus du prix du modèle~?
\item Énoncer le \emph{résultat nommé}~: la prime, son coût annuel et le vega
du cap du trésorier.
\item En une phrase~: que procure un cap que ne procure pas un swap~?
\end{enumerate}
\end{problem}

\section{Questions d'entretien}

\begin{interviewq}[$\star$ \iqroles{trader, banque}]\label{iq:rc:vanilla-rates-options:1}
Qu'est-ce qu'un caplet, et sous quelle probabilité son taux forward est-il une martingale~?
\end{interviewq}

\begin{interviewq}[$\star\star$ \iqroles{chercheur}]\label{iq:rc:vanilla-rates-options:2}
Dérivez la relation entre la volatilité normale et la \omterm{def:rc:vanilla-rates-options:lognormal}{volatilité lognormale} à la monnaie.
\end{interviewq}

\begin{interviewq}[$\star\star$ \iqroles{chercheur, développeur}]\label{iq:rc:vanilla-rates-options:3}
Comment strippez-vous les \omterm{def:rc:vanilla-rates-options:flat}{volatilités des caplets} à partir des cotations de caps, et que peut-il mal se passer~?
\end{interviewq}

\begin{interviewq}[$\star\star$ \iqroles{trader}]\label{iq:rc:vanilla-rates-options:4}
Pourquoi un desk préférerait-il gérer le risque en volatilité normale~?
\end{interviewq}

\begin{interviewq}[$\star\star\star$ \iqroles{chercheur, banque}]\label{iq:rc:vanilla-rates-options:5}
Quel est le problème posé par le pricing d'une \omterm{def:rc:vanilla-rates-options:cash}{swaption à règlement en espèces} au taux au pair avec
la formule de Black et l'\omterm{def:rc:vanilla-rates-options:cash}{annuité cash}, et comment le marché l'a-t-il corrigé~?
\end{interviewq}

\begin{interviewq}[$\star\star\star$ \iqroles{développeur, risque}]\label{iq:rc:vanilla-rates-options:6}
Votre pricer de caplets renvoie NaN pour certaines transactions après une baisse des taux. Pourquoi, et que
modifiez-vous~?
\end{interviewq}
